% Part: first-order-logic
% Chapter: tableaux
% Section: identity

\documentclass[../../../include/open-logic-section]{subfiles}

\begin{document}

\olfileid{fol}{tab}{ide}

\olsection{\usetoken{P}{tableau} with \usetoken{S}{identity}}

带等词的分析表需要额外的推理规则。
关于 $\eq$ 的规则如下（$t$、$t_1$ 和 $t_2$ 均为闭项）：

\begin{defish}
\AxiomC{}
\RightLabel{$\eq$}
\UnaryInfC{\sFmla{\True}{\eq[t][t]}}
\DisplayProof
\hfill
\AxiomC{\sFmla{\True}{\eq[t_1][t_2]}}
\noLine
\UnaryInfC{\sFmla{\True}{!A(t_1)}}
\RightLabel{$\TRule{\True}{\eq}$}
\UnaryInfC{\sFmla{\True}{!A(t_2)}}
\DisplayProof
\hfill
\AxiomC{\sFmla{\True}{\eq[t_1][t_2]}}
\noLine
\UnaryInfC{\sFmla{\False}{!A(t_1)}}
\RightLabel{$\TRule{\False}{\eq}$}
\UnaryInfC{\sFmla{\False}{!A(t_2)}}
\DisplayProof
\end{defish}

与所有其他规则不同，$\TRule{\True}{\eq}$ 和
$\TRule{\False}{\eq}$ 要求分支上\emph{已经出现}两\emph{个}带号公式，
即 $\sFmla{\True}{\eq[t_1][t_2]}$ 与 $\sFmla{S}{!A(t_1)}$ 两者。

\begin{ex}
如果 $s$ 和 $t$ 是闭项，那么 $\eq[s][t], !A(s)
\Proves !A(t)$：
\begin{oltableau}
  [\sFmla{\False}{\formula{A}(t)}, just = \TAss
    [\sFmla{\True}{\eq[s][t]}, just = \TAss
      [\sFmla{\True}{\formula{A}(s)}, just = \TAss
        [\sFmla{\True}{\formula{A}(t)}, just={\TRule{\True}{\eq}[2, 3]}, close]
      ]
    ]
  ]
\end{oltableau}
这或许就是我们熟悉的同一物的可替换性原理，或称莱布尼茨律。

分析表可证明 $\eq$ 是对称的，即 $\eq[s_1][s_2]
\Proves \eq[s_2][s_1]$：
\begin{oltableau}
  [\sFmla{\False}{\eq[s_2][s_1]}, just = \TAss
    [\sFmla{\True}{\eq[s_1][s_2]}, just = \TAss
      [\sFmla{\True}{\eq[s_1][s_1]}, just = {$\eq$}
        [\sFmla{\True}{\eq[s_2][s_1]}, just = {\TRule{\True}{\eq}[2, 3]}, close]
      ]
    ]
  ]
\end{oltableau}
这里，第 $2$ 行是 $\TRule{\True}{\eq}$ 的第一个前提
公式~$\sFmla{\True}{\eq[s_1][s_2]}$。
第 $3$ 行是第二个前提公式，其形式为 $\sFmla{\True}{!A(s_2)}$——不妨
将 $!A(x)$ 视作
$\eq[x][s_1]$，于是 $!A(s_1)$ 即 $\eq[s_1][s_1]$，而 $!A(s_2)$ 即 $\eq[s_2][s_1]$。

分析表也能证明 $\eq$ 是传递的，即 $\eq[s_1][s_2],
\eq[s_2][s_3] \Proves \eq[s_1][s_3]$：
\begin{oltableau}
  [\sFmla{\False}{\eq[s_1][s_3]}, just = \TAss
    [\sFmla{\True}{\eq[s_1][s_2]}, just = \TAss
      [\sFmla{\True}{\eq[s_2][s_3]}, just = \TAss
        [\sFmla{\True}{\eq[s_1][s_3]}, just = {\TRule{\True}{\eq}[3, 2]}, close]
      ]
    ]
  ]
\end{oltableau}
在这个分析表中，$\TRule{\True}{\eq}$ 的第一个前提公式是第 $3$ 行，即 $\sFmla{\True}{\eq[s_2][s_3]}$
（其中 $s_2$ 充当 $t_1$，$s_3$ 充当 $t_2$）。第二个前提公式
形如~$\sFmla{\True}{!A(s_2)}$，即第 $2$ 行。
这里，不妨将 $!A(x)$ 视作 $\eq[s_1][x]$；这样，$!A(s_2)$ 就成为
$\eq[t_1][t_2]$（即第 $2$ 行），而 $!A(s_3)$ 则成为结论中的
公式~$\eq[s_1][s_3]$。
\end{ex}

\begin{prob}
为以下各式构造闭分析表：
\begin{enumerate}
\item $\sFmla{\False}{\lforall[x][\lforall[y][((x = y \land !A(x))
      \lif !A(y))]]}$
\item $\sFmla{\False}{\lexists[x][(!A(x) \land
    \lforall[y][(!A(y) \lif y = x)])]}$,\\
  $\sFmla{\True}{\lexists[x][!A(x)] \land
    \lforall[y][\lforall[z][((!A(y) \land !A(z)) \lif y = z)]]}$
\end{enumerate}
\end{prob}

\end{document}